\section{\raggedright A graded carrier and its interpretation}
Let $H$ be a real Hilbert space and let $J:H\to H$ be an orthogonal
involution. Since $J^{-1}=J=J^*$, the operators
\begin{equation}
 C=\frac{I+J}{2},\qquad P_-=\frac{I-J}{2}
\end{equation}
are complementary orthogonal projections. Indeed,
$C^2=(I+2J+J^2)/4=C$, and similarly $P_-^2=P_-$, while
$CP_-=0$ and $C+P_-=I$. We refer to $C$ as the HC compensator.
It selects the $J$-even part of a state; $P_-$ selects the odd part.

This grading is the minimal structure needed for the subsequent tests.
It does not identify either range with a particle species. In the
Companion, a biquaternionic carrier and extensions involving
$\mathbb C\otimes\mathbb H\otimes\mathbb O$ motivate proposed
particle and frame assignments. Division algebras provide a mathematical
setting for such representation questions \cite{Baez2002}. A positive
Hilbert metric used in a readout calculation must nevertheless be
distinguished from a Lorentzian spacetime metric.

The scalar field also matters. An antiunitary involution on a complex
Hilbert space is generally only real-linear when averaged with the
identity. Here $J$ is explicitly real orthogonal; when complex amplitudes
are used, we mean its declared complex-linear extension. No equivalence
with an inherited antiunitary convention is presumed. Likewise, a change
between the proposed matter, light and tachyon frames requires an
explicit map, not merely a relabelling of components.
\subsection{The algebraic carrier and its two conjugations}
A concrete carrier is needed before the word ``sector'' acquires a
mathematical meaning. Write a biquaternion as
\begin{equation}
 q=q_0+z_1e_1+z_2e_2+z_3e_3,\qquad q_0,z_k\in\mathbb C,
 \qquad e_ie_j=-\delta_{ij}+\epsilon_{ijk}e_k.
\end{equation}
The scalar imaginary unit commutes with the quaternion units. This is
different from choosing a quaternion direction as an internal complex
structure. In particular, a statement formulated with the commuting
scalar unit cannot automatically be transferred to left multiplication
by one of the $e_k$. The representation and the relevant adjoint must
be specified for each transfer.

There are also two distinct conjugations. Define quaternion conjugation
by $\widetilde q=q_0-\sum_kz_ke_k$, leaving the complex coefficients
unchanged, and coefficient conjugation by
$\overline q=\overline{q_0}+\sum_k\overline{z_k}e_k$. Then
\begin{equation}
 q\widetilde q=q_0^2+z_1^2+z_2^2+z_3^2,
\end{equation}
whereas the positive coefficient norm is
\begin{equation}
 \|q\|_E^2=|q_0|^2+|z_1|^2+|z_2|^2+|z_3|^2.
\end{equation}
The first expression is a complex quadratic form; the second is a
positive Hermitian norm. Using the same word ``norm'' for both can
conceal a change of mathematical problem. A positivity argument about
principal mass roots uses a positive comparison. It cannot use a
complex quadratic invariant as though it were that comparison.

One familiar realization sends $e_k$ to $-i\sigma_k$, with the Pauli
matrices $\sigma_k$. The biquaternion becomes a complex two-by-two
matrix, and its determinant is $q_0^2+\sum_kz_k^2$. On the Hermitian
slice $X=tI+\boldsymbol x\cdot\boldsymbol\sigma$ this gives
\begin{equation}
 \det X=t^2-|\boldsymbol x|^2,\qquad
 X\longmapsto AXA^*,\quad A\in SL(2,\mathbb C).
\end{equation}
The determinant is preserved. With the $(-,+,+,+)$ spacetime
convention used in the geometric comparisons, the Lorentz quadratic
form on this slice is $-\det X$. This conventional sign does not
turn the coefficient norm into a Lorentz-invariant Hilbert norm.
For example, $A=\operatorname{diag}(e^{\chi/2},e^{-\chi/2})$
is a boost representative and is not a unitary matrix when
$\chi\ne0$. A similarity, a physical observer transformation and
a unitary change of a positive comparison basis must consequently be
tracked as different operations.

\subsection{What the sign of a field invariant actually distinguishes}
An electromagnetic comparison gives a useful, independently interpretable
instance of the complex vector part. In units with $c=1$, put
$z=\boldsymbol E+i\boldsymbol B$. Then
\begin{equation}
 s=z\cdot z=|\boldsymbol E|^2-|\boldsymbol B|^2
                 +2i\boldsymbol E\cdot\boldsymbol B.
\end{equation}
This identity distinguishes the two real electromagnetic invariants.
Only on a real-$s$ slice, for example
$\boldsymbol E\cdot\boldsymbol B=0$, is a classification by
$s>0$, $s<0$ and $s=0$ directly meaningful. The first two cases are
electric- and magnetic-dominated respectively; the last is a null
invariant when the field is nonzero. Extending a sign label to a
general complex $s$ would require another declared rule.

These algebraic categories are useful for the Companion's proposed
matter, tachyon and photon readouts, but their physical identification
requires dynamics and observables. A null electromagnetic invariant,
for instance, does not by itself specify a quantum photon state.
Nor does a common carrier identify different invariant strata as the
same physical state. The distinction is operational: one must specify
which preparations, transformations and measurements are being compared.

The HC construction is formulated on a positive Hilbert realization
precisely so that its projection identities can be tested independently
of these interpretations. In such a realization every $f$ decomposes
uniquely as
\begin{equation}
 f=f_++f_-,\qquad f_+=Cf,\quad f_-=P_-f,\qquad
 \|f\|^2=\|f_+\|^2+\|f_-\|^2.
\end{equation}
For a normalized state its involution expectation is
\begin{equation}
 \langle f,Jf\rangle=\|f_+\|^2-\|f_-\|^2.
\end{equation}
Zero expectation means equal weights in these two orthogonal sectors.
It is not a fixed-point condition: a $J$-fixed unit vector has expectation
one. This difference is decisive when a symmetry axis or a critical
line is proposed as an explanation of mass-root balance.

\subsection{Realification and the meaning of an adjoint}
Suppose instead that an inherited $J$ is antiunitary on a complex
Hilbert space. It obeys $J(if)=-iJf$, so $(I+J)/2$ is generally not
complex-linear. Regard the same vectors as a real Hilbert space with
inner product $\operatorname{Re}\langle f,g\rangle$. An antiunitary
involution then becomes an orthogonal real-linear involution, and the
projection calculation is valid on that real space. What changes is
the allowed class of operators and subspaces.

One can also complexify a real orthogonal realization and extend its
involution complex-linearly. This produces a different comparison
construction, in which multiplication by $i$ commutes with $J$.
These alternatives cannot be combined in an argument by keeping only
the symbol $J$. Later formulas such as $-iJA$ require the stated
complex-linear extension and a selfadjoint anticommuting $A$ on the
same positive Hilbert space. They are not formulas for an unspecified
antiunitary operation.

The same discipline applies to tensors. A spacetime metric contracts
spacetime indices; an internal Hermitian metric compares internal
amplitudes; a generation metric compares generation coordinates.
Their signatures, dimensions and transformation rules differ. Even
when their numerical matrices happen to coincide in one chart, an
identification must explain how the corresponding spaces and adjoints
are transported. The QAM context records discussed next serve to make
these choices visible at the point where they enter an inference.

\subsection{The physical interface as an additional piece of data}
A physical readout consists of a domain of admitted states together
with a map to observables. For the mass branch this includes a
generation assignment, positivity of principal roots and a dimensional
normalization. For a geometric branch it includes an invertible
coframe or a metric, its differentiation rules and its coupling to
the equations of motion. These requirements are stronger than the
existence of a vector-space isomorphism.

The common-carrier proposal can therefore be investigated in two
directions. Forward construction starts from a represented carrier
and computes its readout. Backward analysis starts from a required
property of the observable and asks which carrier operations could
produce it. The mass-root conditions and the projected-autonomy test
in this paper are backward constraints of this second kind. They do
not discard the common-carrier hypothesis; they specify the extra
mathematical work needed to give it physical content.
